\documentclass[a4paper, serif, 10pt]{moderncv}

%% moderncv
% style and color
\moderncvstyle{banking}
\moderncvcolor{black}

%% page layout/margins
\usepackage[top=2.5cm, bottom=2.5cm, left=1.5cm, right=1.5cm]{geometry}

\nopagenumbers{}

%% Babel config for Dutch language
% ref:
% - https://latex3.github.io/babel/guides/locale-dutch.html
\usepackage[dutch]{babel}

%% fontspec
\usepackage{fontspec}

\IfFontExistsTF{Linux Libertine}{
  \setmainfont[Ligatures={TeX, Common}, Numbers=OldStyle]{Linux Libertine}
}{}
\IfFontExistsTF{Linux Libertine O}{
  \setmainfont[Ligatures={TeX, Common}, Numbers=OldStyle]{Linux Libertine O}
}{}

%% CTeX
% CJK support, used to show CJK characters in the resume
%
% - fontset=none: disable builtin fontset but instead set the CJK font manually
% - heading=false: disable ctex heading
% - punct=kaiming: use kaiming punctuations styles for CJK
% - scheme=plain: use plain scheme, do not override \normalsize font size
% - space=auto: space settings for CJK characters
%
% ref:
% - http://ctan.mirrorcatalogs.com/language/chinese/ctex/ctex.pdf
\usepackage[UTF8, heading=false, punct=kaiming, scheme=plain, space=auto]{ctex}

\IfFontExistsTF{Noto Serif CJK SC}{
  \setCJKmainfont{Noto Serif CJK SC}
}{}
\IfFontExistsTF{Noto Sans CJK SC}{
  \setCJKsansfont{Noto Sans CJK SC}
}{}

%% line spacing
\usepackage{setspace}
\setstretch{1.125}

%% URL styling - use normal text instead of monospace
\urlstyle{same}

% Auto-underline all links
% Defer until \begin{document} so hyperref is already loaded
\AtBeginDocument{%
  \let\oldhref\href
  \renewcommand{\href}[2]{\oldhref{#1}{\underline{#2}}}%
}

%% Patch \httplink and \httpslink to handle full URLs
% The original moderncv macros always prepend http:// or https:// to URLs.
% This causes issues when the URL already has a protocol (e.g., https://example.com)
% resulting in malformed URLs like https://https://example.com.
% This patch checks if the URL already contains "://" and uses it directly if so.
% Note: We use \str_set:Nx (x-expansion) to expand macros like \@homepage first.
\ExplSyntaxOn
\str_new:N \g_yamlresume_url_str

\RenewDocumentCommand{\httpslink}{O{}m}{
  \str_set:Nx \g_yamlresume_url_str {#2}
  \str_if_in:NnTF \g_yamlresume_url_str {://}
    { \url{#2} }
    { \url{https://#2} }
}

\RenewDocumentCommand{\httplink}{O{}m}{
  \str_set:Nx \g_yamlresume_url_str {#2}
  \str_if_in:NnTF \g_yamlresume_url_str {://}
    { \url{#2} }
    { \url{http://#2} }
}
\ExplSyntaxOff

\name{Maarten van der Meer}{}
\title{Senior Projectmanager met 10+ jaar ervaring in digitale transformatie}
\phone[mobile]{+31 6 12345678}
\email{m.vandermeer@example.com}
\homepage{https://www.maartenvandermeer.nl}

\address{Keizersgracht 123, Amsterdam, Noord-Holland, Nederland, 1015 CJ}{}{}

\extrainfo{{\small \faLinkedin}\ \href{https://www.linkedin.com/in/maarten-van-der-meer}{@maarten-van-der-meer} {} {} {} • {} {} {} 
{\small \faMedium}\ \href{https://medium.com/@maartenvandermeer}{@maartenvandermeer}}

\begin{document}

\maketitle

\section{Basis Informatie}

\cvline{}{Ervaren projectmanager met meer dan 10 jaar ervaring in het leiden van complexe, multidisciplinaire projecten. Gespecialiseerd in digitale transformatie, agile methodieken en belanghebbendenmanagement.
%
\begin{itemize}
\item Leidt projecten van initiatie tot oplevering, binnen scope, budget en tijd.
\item Sterk in het verbinden van business en IT, met oog voor mens en resultaat.
\item Gecertificeerd in PRINCE2, PMP en Scrum, en bedreven in Jira, Confluence en MS Project.
\item Coacht en motiveert teams om gezamenlijke doelen te behalen.
\end{itemize}}
\section{Opleidingen}

\cventry{sep 2008–jul 2014}
        {Master, Technische Bestuurskunde, Score: 8.2}
        {Technische Universiteit Delft}
        {\url{https://www.tudelft.nl}}
        {}
        {\begin{itemize}
\item Afgestudeerd met specialisatie in project- en programmamanagement.
\item Ontwikkelde vaardigheden in kwantitatieve analyse, besluitvorming en procesverbetering.
\item Thesis over de invloed van agile werken op publieke projecten.
\end{itemize}
\textbf{Cursussen}: Projectmanagement, Risicomanagement, Operations Research, Financieel Management, Verandermanagement, Systeemkunde}

\cventry{sep 2004–jun 2008}
        {Bachelor, Bestuurskunde, Score: 7.8}
        {Universiteit van Amsterdam}
        {\url{https://www.uva.nl}}
        {}
        {\begin{itemize}
\item Basis gelegd in bestuurskunde, organisatiekunde en beleidsanalyse.
\item Actief in studievereniging en organisatie van congressen.
\end{itemize}
\textbf{Cursussen}: Bestuur en Politiek, Organisatietheorie, Methoden van Onderzoek, Economie}
\section{Werkervaring}

\cventry{mrt 2020–Heden}
        {Senior Projectmanager}
        {ING}
        {\url{https://www.ing.com}}
        {}
        {\begin{itemize}
\item Verantwoordelijk voor een portfolio van digitale transformatieprojecten met een totale waarde van EUR 8 miljoen.
\item Leidt multidisciplinaire teams van 15+ professionals, waaronder ontwikkelaars, data-analisten en zakelijke consultants.
\item Implementeerde een agile manier van werken, wat leidde tot 30\% snellere oplevering en hogere klanttevredenheid.
\item Beheert risico's, budgetten en planningen, en rapporteert aan de stuurgroep.
\item Zorgt voor heldere communicatie en afstemming tussen business, IT en externe leveranciers.
\end{itemize}
\textbf{Trefwoorden}: Agile, Scrum, Digitaal, Belanghebbenden, Risicobeheer}

\cventry{jan 2016–feb 2020}
        {Projectmanager}
        {Philips}
        {\url{https://www.philips.com}}
        {}
        {\begin{itemize}
\item Leidde productontwikkelingsprojecten in de gezondheidszorgtechnologie, van concept tot marktlancering.
\item Werkte samen met R\&D, marketing en regelgevende zaken om naleving en tijdige oplevering te waarborgen.
\item Introduceerde Scrum en verbeterde de voorspelbaarheid van opleveringen met 25\%.
\item Beheerde een projectbudget van EUR 3 miljoen en een team van 10 projectleden.
\end{itemize}
\textbf{Trefwoorden}: Productontwikkeling, Gezondheidszorg, Scrum, Budgetbeheer}

\cventry{sep 2014–dec 2015}
        {Junior Projectmanager}
        {KPN}
        {\url{https://www.kpn.com}}
        {}
        {\begin{itemize}
\item Ondersteunde senior projectmanagers bij telecom-infrastructuurprojecten.
\item Coördineerde planningen, vergaderingen en rapportages voor interne en externe belanghebbenden.
\item Droeg bij aan de uitrol van glasvezelnetwerken in Noord-Holland.
\end{itemize}
\textbf{Trefwoorden}: Telecom, Planning, Rapportage, Glasvezel}
\section{Talen}

\cvline{Nederlands}{Moedertaal of tweetalige vaardigheid}
\cvline{Engels}{Volledige professionele vaardigheid \hfill \textbf{Trefwoorden}: IELTS 8.0, Zakelijk vloeiend}
\cvline{Duits}{Beperkte werkvaardigheid \hfill \textbf{Trefwoorden}: Goethe B1, Basiskennis}
\section{Vaardigheden}

\cvline{Projectmanagement}{Expert \hfill \textbf{Trefwoorden}: Agile, Scrum, PRINCE2, Waterval, Jira, MS Project}
\cvline{Belanghebbendenmanagement}{Geavanceerd \hfill \textbf{Trefwoorden}: Communicatie, Onderhandelen, Rapportage, Workshops, Verandermanagement}
\cvline{Technisch}{Geavanceerd \hfill \textbf{Trefwoorden}: SQL, Python, Power BI, Confluence, Azure DevOps}
\cvline{Leiderschap}{Expert \hfill \textbf{Trefwoorden}: Coaching, Teambuilding, Motiveren, Conflicthantering}
\section{Onderscheidingen}

\cventry{nov 2022}
        {PMI Netherlands Chapter}
        {Projectmanager van het Jaar}
        {}
        {}
        {\begin{itemize}
\item Erkend voor uitmuntend leiderschap en innovatie in projectmanagement.
\item Bijzondere vermelding voor het succesvol opleveren van een complex digitaal transformatieprogramma.
\end{itemize}}

\cventry{jun 2021}
        {ING}
        {Beste Innovatieproject}
        {}
        {}
        {\begin{itemize}
\item Toegekend voor de ontwikkeling van een realtime klantdashboard dat de Net Promoter Score met 15 punten verhoogde.
\end{itemize}}
\section{Certificaten}

\cventry{mei 2018}
        {AXELOS}
        {PRINCE2 Practitioner}
        {\url{https://www.axelos.com/certifications/prince2}}
        {}
        {}

\cventry{sep 2020}
        {Project Management Institute}
        {Projectmanagement Professional (PMP)}
        {\url{https://www.pmi.org/certifications/project-management-pmp}}
        {}
        {}

\cventry{mrt 2017}
        {Scrum Alliance}
        {Gecertificeerd ScrumMaster}
        {\url{https://www.scrumalliance.org/get-certified/scrum-master-track/certified-scrummaster}}
        {}
        {}
\section{Publicaties}

\cventry{apr 2021}
        {Agile projectmatig werken in de publieke sector}
        {Management \& Organisatie}
        {\url{https://www.managementenorganisatie.nl/agile-publieke-sector}}
        {}
        {\begin{itemize}
\item Onderzoek naar de toepassing van agile methodieken binnen overheidsprojecten.
\item Biedt praktische handvatten voor projectmanagers in een politiek-bestuurlijke context.
\end{itemize}}
\section{Projecten}

\cventry{feb 2021–dec 2022}
        {Een organisatiebreed programma om de digitale dienstverlening van ING te moderniseren en te versnellen.}
        {Digitale Transformatie ING}
        {}
        {}
        {\begin{itemize}
\item Leidde de implementatie van een nieuw digitaal klantplatform voor 2 miljoen klanten.
\item Verantwoordelijk voor planning, budget, risicobeheer en afstemming met belanghebbenden.
\item Resultaat: 30\% snellere klantinteracties en 20\% lagere operationele kosten.
\end{itemize}
\textbf{Trefwoorden}: Digitale Transformatie, Klantplatform, Agile}

\cventry{mei 2019–nov 2019}
        {Een data-gedreven dashboard voor het meten en rapporteren van duurzaamheidsprestaties.}
        {Duurzaamheidsmonitor Philips}
        {}
        {}
        {\begin{itemize}
\item Ontwikkelde een monitor om de CO2-uitstoot van productielocaties direct te volgen.
\item Werkte samen met duurzaamheidsteam en data-ingenieurs.
\item Resultaat: inzicht in 15\% reductie van energieverbruik.
\end{itemize}
\textbf{Trefwoorden}: Duurzaamheid, Data, Dashboard}
\section{Interesses}

\cvline{Sport}{Hardlopen, Wielrennen, Zwemmen}
\cvline{Cultuur}{Musea, Literatuur, Reizen}
\section{Vrijwilligerswerk}

\cventry{sep 2017–Heden}
        {Mentor}
        {PMI Netherlands Chapter}
        {\url{https://www.pmi-netherlands-chapter.org}}
        {}
        {\begin{itemize}
\item Begeleidt startende projectmanagers bij hun professionele ontwikkeling.
\item Organiseert maandelijkse netwerkbijeenkomsten en workshops over projectmanagement.
\item Geeft presentaties over agile leiderschap en belanghebbendenmanagement.
\end{itemize}}
    
\end{document}